\section*{Розділ III. Процедури перших виборів делегатів КСУ}
\addcontentsline{toc}{section}{Розділ III. Процедури перших виборів делегатів КСУ}

\subsection*{3.1. Загальна послідовність проведення виборів делегатів КСУ}
\addcontentsline{toc}{subsection}{3.1. Загальна послідовність проведення виборів делегатів КСУ}
    3.1.1. Перші вибори проводяться виключно для обрання делегатів до Конференції студентів Університету (КСУ), після чого скликається перше засідання КСУ.

    3.1.2. Вибори делегатів КСУ проводяться протягом першого місяця після формування ТВК.

    3.1.3. Перше засідання КСУ скликається протягом 2 тижнів після завершення виборів делегатів.
    3.1.4. Після оголошення результатів голосування та скликання першого засідання КСУ формування інших органів студентського самоврядування здійснюється за положеннями, затвердженими КСУ.

\subsection*{3.2. Строки проведення виборчої кампанії}
\addcontentsline{toc}{subsection}{3.2. Строки проведення виборчої кампанії}
    3.2.1. Строк проведення виборчої кампанії для виборів делегатів КСУ становить 10 календарних днів.

    3.2.2. Виборча кампанія включає такі етапи:

        \begin{enumerate}[label=\alph*)]
            \item Оголошення про початок виборів -- 1 день;
            \item Висування та реєстрація кандидатів -- 4 дні;
            \item Агітаційний період -- 3 дні;
            \item День тиші -- 1 день перед голосуванням;
            \item Голосування -- 1 день;
            \item Підрахунок голосів та оголошення результатів -- не пізніше 24 годин з моменту завершення голосування.
        \end{enumerate}

    3.2.3. ТВК має право скоротити строки окремих етапів виборчої кампанії, але не більше ніж на 30\% від встановленого строку, за умови забезпечення достатнього часу для реалізації виборчих прав студентів.

\subsection*{3.3. Виборці та виборчі права}
\addcontentsline{toc}{subsection}{3.3. Виборці та виборчі права}
    3.3.1. Право голосу на перших виборах мають усі студенти Університету, які відповідають вимогам Положення про органи студентського самоврядування щодо участі в ОСС.

    3.3.2. Перші вибори визнаються дійсними за умови участі у голосуванні не менше 1 кандидата та наявності будь-якої кількості виборців. У разі участі у виборах лише одного кандидата, він вважається обраним, якщо за нього проголосувало більше 50\% виборців, що взяли участь у голосуванні.

    3.3.3. У разі відсутності кандидатів, ТВК продовжує етап висування кандидатів на додатковий строк до 7 днів.

\subsection*{3.4. Висування та реєстрація кандидатів}
\addcontentsline{toc}{subsection}{3.4. Висування та реєстрація кандидатів}
    3.4.1. Кандидати на перші вибори висуваються у спосіб, передбачений відповідними положеннями про ОСС, з урахуванням спрощень, встановлених цим Регламентом.

    3.4.2. Для реєстрації кандидата подаються:

        \begin{enumerate}[label=\alph*)]
            \item Заява про згоду балотуватися;
            \item Підтвердження статусу студента Університету (довідка, студентський квиток, перепустка в гуртожиток або інший документ, що засвідчує статус студента);
            \item Програма кандидата (не більше 1 сторінки А4);
        \end{enumerate}

    3.4.3. ТВК розглядає документи кандидатів протягом 2 робочих днів та приймає рішення про реєстрацію або відмову в реєстрації з обґрунтуванням.

    3.4.4. Кандидат має право оскаржити рішення про відмову в реєстрації до ТВК протягом 2 (двох) робочих днів. ТВК зобов'язана розглянути скаргу протягом 2 (двох) робочих днів.

    3.4.5. У разі незгоди з рішенням ТВК за результатами розгляду скарги, кандидат має право звернутися до Загальних зборів студентів Університету протягом 2 (двох) робочих днів з дня отримання рішення ТВК. Рішення Загальних зборів студентів є остаточним та не підлягає оскарженню.

\subsection*{3.5. Агітаційний період}
\addcontentsline{toc}{subsection}{3.5. Агітаційний період}
    3.5.1. Агітаційна діяльність під час перших виборів проводиться у спрощеному форматі та включає:

        \begin{enumerate}[label=\alph*)]
            \item Розміщення програм кандидатів на офіційних інформаційних ресурсах ТВК;
            \item Розповсюдження друкованих матеріалів кандидатів у визначених місцях;
            \item Проведення презентацій кандидатів (за бажанням та можливістю);
            \item Розміщення інформації в соціальних мережах та месенджерах.
        \end{enumerate}

    3.5.2. Заборонено під час агітаційного періоду:

        \begin{enumerate}[label=\alph*)]
            \item Використання адміністративного ресурсу;
            \item Підкуп виборців у будь-якій формі;
            \item Поширення неправдивої інформації про кандидатів;
            \item Агітацію в день голосування та день тиші перед ним.
        \end{enumerate}

    3.5.3. ТВК забезпечує рівні можливості для агітації всіх зареєстрованих кандидатів, включаючи рівний доступ до офіційних інформаційних ресурсів.

    3.5.4. Фінансування агітаційної діяльності здійснюється виключно за рахунок власних коштів кандидатів або добровільних пожертв фізичних осіб у розмірі, що не перевищує встановлені ТВК ліміти.

\subsection*{3.6. Організація голосування}
\addcontentsline{toc}{subsection}{3.6. Організація голосування}
    3.6.1. Голосування на перших виборах може проводитися у таких формах:

        \begin{enumerate}[label=\alph*)]
            \item Очне голосування на виборчих дільницях;
            \item Електронне голосування через офіційні платформи Університету;
            \item Змішане голосування (поєднання очного та електронного).
        \end{enumerate}

    3.6.2. Форму голосування для кожного типу виборів визначає ТВК з урахуванням технічних можливостей та специфіки виборців.

    3.6.3. ТВК створює дільничні виборчі комісії для організації голосування на виборчих дільницях. До складу дільничних комісій можуть входити студенти, які не є кандидатами на відповідних виборах.

    3.6.4. Голосування проводиться протягом одного дня з 9:00 до 18:00, якщо ТВК не встановить інший час з урахуванням специфіки відповідних виборів.

    3.6.5. За порядком на виборчих дільницях слідкують члени дільничних комісій. Кандидати мають право призначати довірених осіб для спостереження за ходом голосування.

\subsection*{3.7. Підрахунок голосів та оголошення результатів}
\addcontentsline{toc}{subsection}{3.7. Підрахунок голосів та оголошення результатів}
    3.7.1. Підрахунок голосів здійснюється відкрито за присутності членів дільничних комісій, довірених осіб кандидатів та спостерігачів.

    3.7.2. Результати підрахунку голосів на кожній дільниці оформлюються протоколом, який підписується всіма присутніми членами дільничної комісії.

    3.7.3. ТВК узагальнює результати голосування та не пізніше наступного дня після голосування оголошує остаточні результати виборів.

    3.7.4. Обраними вважаються кандидати, які отримали найбільшу кількість голосів виборців у межах квоти відповідного виборчого округу.

    3.7.5. У разі отримання однакової кількості голосів двома або більше кандидатами, ТВК призначає повторне голосування між цими кандидатами протягом 5 днів.

    3.7.6. Результати виборів оприлюднюються на всіх офіційних інформаційних ресурсах ТВК та Університету не пізніше ніж через 24 години після їх затвердження.